\begin{table}[htbp]
\centering
\caption{Covariate Balance at the Threshold}
\label{tab:covariate_balance}
\begin{tabular}{l ccccc}
\toprule
Covariate & RDD Estimate & Robust SE & $p$-value & Bandwidth & Eff.\ $N$ \\
\midrule
Turnout (\%)                   & -2.214 & 1.874 & 0.237 & 11.19 & 1,273 \\
No.\ Candidates                & 0.917 & 0.673 & 0.173 & 11.81 & 1,310 \\
Valid Votes                    & -920.874 & 5240.296 & 0.861 & 13.07 & 1,417 \\
Total Electors                 & 4215.158 & 9438.148 & 0.655 & 13.46 & 1,436 \\
Log Population (2001)          & 0.067 & 0.070 & 0.338 & 11.83 & 1,159 \\
Literacy Rate (2001)           & -0.002 & 0.021 & 0.932 & 10.20 & 1,056 \\
SC Share (2001)                & -0.013 & 0.012 & 0.277 & 14.61 & 1,339 \\
ST Share (2001)                & -0.016 & 0.022 & 0.483 & 11.40 & 1,134 \\
Pre-Election NL                & 199.157 & 530.400 & 0.707 & 10.70 & 1,230 \\
\bottomrule
\end{tabular}
\begin{minipage}{0.92\textwidth}
\vspace{4pt}
\footnotesize
\textit{Notes:} Each row reports a separate RDD regression of the covariate on the
criminal-candidate vote margin, using local linear estimation with MSE-optimal bandwidth.
Bias-corrected estimates with robust standard errors.
No significant differences indicate the RDD design is valid.
$^{***}$ $p<0.01$, $^{**}$ $p<0.05$, $^{*}$ $p<0.1$.
\end{minipage}
\end{table}
